\section{Partie théorique} \label{sec:theorie}

\subsection{Oscillations libres d’un pendule de Pohl} \label{subsec:osc_libres}

Le pendule de Pohl est un oscillateur mécanique rotatif composé d'un disque couplé à un ressort
spiral. La coordonnée généralisée du système est l’angle de torsion $\theta(t)$.
En appliquant le théorème du moment cinétique au disque, on obtient :

\begin{equation}
\frac{dL}{dt} = \sum M_{\mathrm{ext}},
\qquad L = J\,\dot{\theta},
\label{eq:TMC}
\end{equation}

où $J$ est le moment d’inertie du balancier. Deux moments mécaniques agissent sur le système :
\begin{equation}
M_r = -C\,\theta,
\qquad
M_f = -b\,\dot{\theta},
\label{eq:MOMENTS}
\end{equation}
où $C$ est la constante de torsion, et $b$ le moment de frottement visqueux.  
En combinant ces contributions, l'équation du mouvement devient :

\begin{equation}
J\ddot{\theta} + b\dot{\theta} + C\theta = 0.
\label{eq:edo}
\end{equation}

Dont on définit la pulsation propre non amortie $\omega_0$ et le coefficient d'amortissement $\lambda$ tels que :
\begin{equation}
\omega_0 = \sqrt{\frac{C}{J}},
\qquad
\lambda = \frac{b}{2J}.
\label{eq:lambda} % et omega 0
\end{equation}

L’équation différentielle se met sous la forme caractéristique :
\begin{equation}
\ddot{\theta} + 2\lambda \dot{\theta} + \omega_0^2 \theta = 0.
\label{eq:cara}
\end{equation}


